%

\subsection{Johann Lambert} % lang-pl
\label{sekcja_lambert}

Trzydzieści trzy lata po książeczce Saccheriego, Johann Heinrich Lambert napisze podobną pracę badawczą \emph{,,Die Theorie der Parallellinien''}\footnote{Teoria prostych równoległych.}, chociaż ta publikacji doczeka się dopiero jedenaście lat po śmierci autora.
Lambert weźmie połowę czworokątu Saccheriego, która ma trzy kąty proste (ćwiczenia 8, 9 u Evesa \cite[s. 289]{eves1_1972}) i zastanowi się, jaką miarę może mieć czwarty kąt:

\begin{definition}
[czworokąt Lamberta] % lang-pl
	Czworokąt o trzech kątach prostych. % lang-pl
\index{czworokąt!Lamberta} % lang-pl
\end{definition}

Zajdzie dalej niż Saccheri: pokaże, że suma kątów trójkąta jest mniejsza, równa lub większa dwóm kątom prostym oraz że różnica między tą sumą oraz dwoma kątami prostymi jest proporcjonalna do pola trójkąta (w przypadku ostrym i rozwartym).

Wysnuje przypuszczenie, że skoro hipoteza kąta rozwartego przypomina geometrię sferyczną, to hipoteza kąta ostrego może odpowiadać geometrii na sferze o urojonym promieniu.
Być może brak jednoznacznego rozstrzygnięcia w tym przypadku będzie powodem, dlaczego jego wyniki zostaną opublikowane dopiero po śmierci.

O Lambercie pisze Eves \cite[s. 285, 286]{eves1_1972}, Hartshorne w ćwiczeniu \cite[s. 316]{hartshorne_2000}.

%